% TOP_PROBLEM: {"id": 6700013, "title": "Scalar Curvature Question [?13]: [∗] no closed aspherical15 manifold admits a metric withSc > 0", "category_id": 6, "category": {"id": 6, "name": "geometry", "display_name": "Geometry", "description": "Euclidean and non-Euclidean geometry, geometric structures.", "slug": "geometry", "order_index": 6, "created_at": "2026-07-31T15:26:25.671Z"}, "status": "partially_solved"}
% FIRST_POSTED: 2026-10-06
% ATTEMPT_STATUS: solved
% OPENAI_CLAIM: {"schema_version":1,"source_repo":"https://github.com/openai/math","source_commit":"adc7f1241b42e322a6451854ab7e4b4c146bf78a","release_date":"2026-10-06","match":"full","resolution":"proved","independently_reviewed":false,"second_pass":"agreed","families":[{"family":"335","title":"Gromov’s integral scalar-curvature bound for simplicial volume","subject":"Differential geometry","manuscripts":[{"title":"Positive scalar curvature forces rational inessentiality","pdf_path":"preprints/Positive-scalar-curvature-forces-rational-inessentiality-September-23-2026/paper.pdf","readme_path":"preprints/Positive-scalar-curvature-forces-rational-inessentiality-September-23-2026/README.md","date":"2026-09-23","theorem_ref":"thm:main (sections/01-introduction.tex:22); cor:aspherical (sections/01-introduction.tex:184)","statement_sources":[{"path":"preprints/Positive-scalar-curvature-forces-rational-inessentiality-September-23-2026/build/sections/01-introduction.tex","line":22,"label":"thm:main"},{"path":"preprints/Positive-scalar-curvature-forces-rational-inessentiality-September-23-2026/build/sections/01-introduction.tex","line":184,"label":"cor:aspherical"}]}],"lean":null,"reasoning_trace":null}]}
\documentclass[11pt]{article}
\usepackage[margin=25mm]{geometry}
\usepackage{fontspec}
\setmainfont{Latin Modern Roman}
\usepackage{amsmath,amssymb,mathtools}
\usepackage{unicode-math}
\setmathfont{Latin Modern Math}
\usepackage{enumitem}
\usepackage{xurl,hyperref,bookmark}
\hypersetup{colorlinks=true,urlcolor=blue}
\setcounter{secnumdepth}{0}
\setlength{\parindent}{0pt}
\setlength{\parskip}{5pt}
\setlength{\emergencystretch}{3em}
\begin{document}
\section*{Scalar Curvature Question [?13]: [\(\ast\)] no closed aspherical15 manifold admits a metric withSc > 0}

\subsection{OpenAI's claim}
\textit{Positive scalar curvature and inessentiality}: OpenAI claims that positive scalar curvature forces the rational fundamental class to vanish under the classifying map, without spin or group restrictions. Consequences include the aspherical obstruction, flatness rigidity at nonnegative scalar curvature, and zero real simplicial volume in the nonnegative case.

\subsection{Mathematical statement of the claimed result}
\paragraph{Mathematical claim 1.} For a closed connected oriented smooth $n$-manifold, $n\ge2$, admitting a metric with $\operatorname{Scal}>0$, the classifying map $c_M:M\to B\pi_1(M)$ is claimed to satisfy \[(c_M)_*[M]=0\in H_n(B\pi_1(M);\mathbb Q).\] Its explicit aspherical corollary rules out positive scalar curvature on every closed smooth aspherical manifold, including nonorientable ones by passage to an orientable cover.

\subsection{Outline of the argument}
\textit{Positive scalar curvature and inessentiality}: An essential rational class is detected on an auxiliary cover at a finite word-metric scale. Compactly supported cohomology and relative duality produce a finite-dimensional comparison with a nonzero local degree. Transverse realization and scalar-curvature reduction contradict that degree; the nonnegative aspherical rigidity statement then uses splitting after the positive-curvature obstruction is established.

\subsection{Scope relative to this problem}
\textit{Positive scalar curvature and inessentiality}: The explicit aspherical corollary rules out positive scalar curvature for every closed smooth aspherical manifold, including nonorientable ones; this is the entire recorded obstruction.

With the hypotheses and conventions above, the claimed result addresses the whole mathematical assertion on this page.

\subsection{Formal verification}
\textit{Positive scalar curvature and inessentiality}: The pinned catalogue lists no Lean formalization for this result.

This site has not run the Lean development. These links describe OpenAI’s selected formal statements and should be read with their stated scope.

\subsection{Status caveat}
OpenAI reports that an internal model produced the manuscripts. Their arguments have not been independently verified as proofs by this site. Any portion outside the documented formal scope may contain mistakes, and the catalogue label reports the announced mathematical claim.

\subsection{Sources}
\begin{itemize}
\item OpenAI, \emph{Positive scalar curvature forces rational inessentiality}, 2026-09-23; citation key \texttt{OAI:Positive-scalar-curvature-forces-rational-inessentiality-September-23-2026}. \href{https://github.com/openai/math/blob/adc7f1241b42e322a6451854ab7e4b4c146bf78a/preprints/Positive-scalar-curvature-forces-rational-inessentiality-September-23-2026/paper.pdf}{PDF (GitHub)}; \href{https://github.com/openai/math/raw/adc7f1241b42e322a6451854ab7e4b4c146bf78a/preprints/Positive-scalar-curvature-forces-rational-inessentiality-September-23-2026/paper.pdf}{PDF (direct)}; \href{https://github.com/openai/math/blob/adc7f1241b42e322a6451854ab7e4b4c146bf78a/preprints/Positive-scalar-curvature-forces-rational-inessentiality-September-23-2026/README.md}{Manuscript page and citation}; \href{https://github.com/openai/math/blob/main/preprints/Positive-scalar-curvature-forces-rational-inessentiality-September-23-2026/paper.pdf}{Latest version}; Statement source: \href{https://github.com/openai/math/blob/adc7f1241b42e322a6451854ab7e4b4c146bf78a/preprints/Positive-scalar-curvature-forces-rational-inessentiality-September-23-2026/build/sections/01-introduction.tex\#L22}{thm:main}; Statement source: \href{https://github.com/openai/math/blob/adc7f1241b42e322a6451854ab7e4b4c146bf78a/preprints/Positive-scalar-curvature-forces-rational-inessentiality-September-23-2026/build/sections/01-introduction.tex\#L184}{cor:aspherical}.
\item \href{https://github.com/openai/math/blob/adc7f1241b42e322a6451854ab7e4b4c146bf78a/CONTENTS.md}{OpenAI catalogue entry}.
\end{itemize}
\end{document}
